% !TeX program = pdflatex
\documentclass[12pt,a4paper]{article}

\newcommand{\notelanguage}{finnish}
% Shared preamble for course notes and exercise sets.

\usepackage[T1]{fontenc}
\usepackage{lmodern}
\usepackage[english]{babel}

\usepackage[a4paper,margin=30mm]{geometry}
\usepackage{microtype}
\usepackage{graphicx}
\usepackage{enumitem}

\usepackage{amsmath}
\usepackage{amssymb}
\usepackage{amsthm}
\usepackage{mathtools}
\usepackage{aliascnt}

\usepackage[hidelinks,pdfusetitle]{hyperref}

\numberwithin{equation}{section}
\setlist{itemsep=0.25em,topsep=0.5em}

\theoremstyle{plain}
\newtheorem{theorem}{Theorem}[section]
\newaliascnt{lemma}{theorem}
\newtheorem{lemma}[lemma]{Lemma}
\aliascntresetthe{lemma}
\newaliascnt{proposition}{theorem}
\newtheorem{proposition}[proposition]{Proposition}
\aliascntresetthe{proposition}
\newaliascnt{corollary}{theorem}
\newtheorem{corollary}[corollary]{Corollary}
\aliascntresetthe{corollary}

\theoremstyle{definition}
\newaliascnt{definition}{theorem}
\newtheorem{definition}[definition]{Definition}
\aliascntresetthe{definition}
\newaliascnt{example}{theorem}
\newtheorem{example}[example]{Example}
\aliascntresetthe{example}
\newaliascnt{problem}{theorem}
\newtheorem{problem}[problem]{Problem}
\aliascntresetthe{problem}

\theoremstyle{remark}
\newaliascnt{remark}{theorem}
\newtheorem{remark}[remark]{Remark}
\aliascntresetthe{remark}

\usepackage[nameinlink,noabbrev]{cleveref}

\crefname{theorem}{theorem}{theorems}
\crefname{lemma}{lemma}{lemmas}
\crefname{proposition}{proposition}{propositions}
\crefname{corollary}{corollary}{corollaries}
\crefname{definition}{definition}{definitions}
\crefname{example}{example}{examples}
\crefname{problem}{problem}{problems}
\crefname{remark}{remark}{remarks}

\DeclareMathOperator{\im}{im}
\DeclareMathOperator{\rank}{rank}
\DeclarePairedDelimiter{\abs}{\lvert}{\rvert}
\DeclarePairedDelimiter{\norm}{\lVert}{\rVert}
\DeclarePairedDelimiter{\set}{\lbrace}{\rbrace}


\usepackage{tasks}
\usepackage{tikz}
\usepackage{caption}

\settasks{
  label = (\alph*),
  label-width = 2em,
  after-item-skip = 0pt
}

\title{Johdatus yliopistomatematiikkaan\\Harjoituksen 2 ratkaisut}
\date{}
\author{}

\begin{document}

\maketitle

\noindent

\begin{exercise}
  Kirjoita käyttäen matemaattista notaatiota (joukko-opillisia ja loogisia
  merkkejä)
\begin{enumerate}[label=\alph*), nosep]
  \item luku \(\sqrt{2}\) ei kuulu rationaalilukuihin,
  \item lukujen \(-7\), \(0\), \(51\) ja miljoona muodostama joukko on
    kokonaislukujen osajoukko,
  \item joukko \(C\) on joukkojen \(A\) ja \(B\) leikkauksen komplementin
    aito osajoukko,
  \item niiden reaalilukujen \(x\) joukko, joiden neliö on suurempi kuin
    \(1{,}2\) ja pienempi kuin \(3{,}6\),
  \item lukujen \(a\) ja \(b\) irrationaalisuudesta ei seuraa summan
    \(a + b\) irrationaalisuus,
  \item joukkojen \(A\) ja \(B\) yhdiste sisältää joukkojen \(C\) ja \(D\)
    leukkauksen,
  \item jokaista nollasta poikkeavaa kokonaislukua \(x\) kohti on olemassa
    sellainen luku \(y\), että \(xy = 2\),
  \item \(\RR\):n suljetun välin kolmesta viiteen ja avoimen välin neljästä
    kuuteen joukkoerotus on suljettu väli kolmesta neljään.
\end{enumerate}
\end{exercise}

\begin{solution}
\leavevmode\par
\begin{enumerate}[label=\alph*), nosep]
\item \(\sqrt{2} \notin \QQ\)
\item \(\set{-7, 0, 51, 10^6} \subseteq \ZZ\)
\item \(C \subset \complement{(A \cap B)}\)
\item \(\set{x \in \RR \colon 1{,}2 < x^2 < 3{,}6}\)
\item \(\bigl((a \in \RR \setminus \QQ) \land (b \in \RR \setminus \QQ)\bigr)
  \nRightarrow a + b \in \RR \setminus \QQ\)
\item \(C \cap D \subseteq A \cup B\)
\item \(\forall x \in \ZZ \setminus \set{0} \; \exists y \in \RR
  \colon xy = 2\)
\item \([3, 5] \setminus \left]4, 6\right[ = [3, 4]\)
\end{enumerate}
\end{solution}

\begin{exercise}
Olkoon \(X = \set{a, b, c, d, e}\) perusjoukko sekä
\(A = \set{a, b, c}\) ja \(B = \set{b, c, d}\).
Määritä
\begin{enumerate}[label=\alph*), nosep]
\item \(A \cup B\) ja \(A \cap B\),
\item \(\complement{A}\) ja \(\complement{B}\),
\item \(A\):n ja \(B\):n \emph{symmetrinen erotus} \(A \mathbin{\Delta} B = (A
\setminus B) \cup (B \setminus A)\),
\item \(B\):n \emph{potenssijoukko} \(\powerset{B} = \set{C : C
\subseteq B}\).
\end{enumerate}
Havainnollista joukkoja Vennin diagrammeilla.
\end{exercise}

\begin{solution}
\leavevmode\par
\begin{enumerate}[label=\alph*), nosep]
\item \(A \cup B = \set{a, b, c, d}\) ja \(A \cap B = \set{b, c}\)
\item \(\complement{A} = \set{d, e}\) ja \(\complement{B} = \set{a, e}\)
\item \(A \mathbin{\Delta} B = \set{a} \cup \set{d} = \set{a, d}\)
\item \(\powerset{B} = \set{\emptyset, \set{b}, \set{c}, \set{d},
\set{b, c}, \set{b, d}, \set{c, d}, \set{b, c, d}}\)
\end{enumerate}

\begin{center}
\begin{minipage}{\textwidth}
\centering
\begin{tikzpicture}[scale=0.9, every node/.style={font=\small}]
% Union
\begin{scope}
  \fill[blue!25] (1.4,1.2) circle (0.85);
  \fill[blue!25] (2.2,1.2) circle (0.85);
  \draw (0,0) rectangle (3.6,2.7);
  \draw (1.4,1.2) circle (0.85);
  \draw (2.2,1.2) circle (0.85);
  \node at (0.2,2.48) {$X$};
  \node at (1.4,2.4) {$A$};
  \node at (2.2,2.4) {$B$};
  \node at (1.0,1.2) {$a$};
  \node at (1.8,1.45) {$b$};
  \node at (1.8,0.95) {$c$};
  \node at (2.6,1.2) {$d$};
  \node at (3.25,0.3) {$e$};
  \node[below] at (1.8,0) {$A \cup B$};
\end{scope}

% Intersection
\begin{scope}[xshift=5cm]
  \begin{scope}
    \clip (1.4,1.2) circle (0.85);
    \fill[blue!25] (2.2,1.2) circle (0.85);
  \end{scope}
  \draw (0,0) rectangle (3.6,2.7);
  \draw (1.4,1.2) circle (0.85);
  \draw (2.2,1.2) circle (0.85);
  \node at (0.2,2.48) {$X$};
  \node at (1.4,2.4) {$A$};
  \node at (2.2,2.4) {$B$};
  \node at (1.0,1.2) {$a$};
  \node at (1.8,1.45) {$b$};
  \node at (1.8,0.95) {$c$};
  \node at (2.6,1.2) {$d$};
  \node at (3.25,0.3) {$e$};
  \node[below] at (1.8,0) {$A \cap B$};
\end{scope}

% Complement of A
\begin{scope}[yshift=-3.4cm]
  \fill[blue!25] (0,0) rectangle (3.6,2.7);
  \fill[white] (1.4,1.2) circle (0.85);
  \draw (0,0) rectangle (3.6,2.7);
  \draw (1.4,1.2) circle (0.85);
  \draw (2.2,1.2) circle (0.85);
  \node at (0.2,2.48) {$X$};
  \node at (1.4,2.4) {$A$};
  \node at (2.2,2.4) {$B$};
  \node at (1.0,1.2) {$a$};
  \node at (1.8,1.45) {$b$};
  \node at (1.8,0.95) {$c$};
  \node at (2.6,1.2) {$d$};
  \node at (3.25,0.3) {$e$};
  \node[below] at (1.8,0) {$\complement{A}$};
\end{scope}

% Complement of B
\begin{scope}[xshift=5cm,yshift=-3.4cm]
  \fill[blue!25] (0,0) rectangle (3.6,2.7);
  \fill[white] (2.2,1.2) circle (0.85);
  \draw (0,0) rectangle (3.6,2.7);
  \draw (1.4,1.2) circle (0.85);
  \draw (2.2,1.2) circle (0.85);
  \node at (0.2,2.48) {$X$};
  \node at (1.4,2.4) {$A$};
  \node at (2.2,2.4) {$B$};
  \node at (1.0,1.2) {$a$};
  \node at (1.8,1.45) {$b$};
  \node at (1.8,0.95) {$c$};
  \node at (2.6,1.2) {$d$};
  \node at (3.25,0.3) {$e$};
  \node[below] at (1.8,0) {$\complement{B}$};
\end{scope}

% Symmetric difference
\begin{scope}[xshift=2.5cm,yshift=-6.8cm]
  \fill[blue!25] (1.4,1.2) circle (0.85);
  \fill[blue!25] (2.2,1.2) circle (0.85);
  \begin{scope}
    \clip (1.4,1.2) circle (0.85);
    \fill[white] (2.2,1.2) circle (0.85);
  \end{scope}
  \draw (0,0) rectangle (3.6,2.7);
  \draw (1.4,1.2) circle (0.85);
  \draw (2.2,1.2) circle (0.85);
  \node at (0.2,2.48) {$X$};
  \node at (1.4,2.4) {$A$};
  \node at (2.2,2.4) {$B$};
  \node at (1.0,1.2) {$a$};
  \node at (1.8,1.45) {$b$};
  \node at (1.8,0.95) {$c$};
  \node at (2.6,1.2) {$d$};
  \node at (3.25,0.3) {$e$};
  \node[below] at (1.8,0) {$A \mathbin{\Delta} B$};
\end{scope}
\end{tikzpicture}
\captionof{figure}{Kohtien (a)--(c) joukkojen visualisoinnit Vennin
diagrammeina.}
\end{minipage}
\end{center}

Kohdassa (d) valitaan perusjoukoksi \(\powerset{B}\). Piirretään sen sisään
kolme joukkoa, joista ensimmäiseen kuuluvat \(b\):n sisältävät osajoukot,
toiseen \(c\):n sisältävät osajoukot ja kolmanteen \(d\):n sisältävät
osajoukot. Tällöin Vennin diagrammin kahdeksan aluetta vastaavat täsmälleen
joukon \(B\) kahdeksaa osajoukkoa.

\begin{center}
\begin{minipage}{\textwidth}
\centering
\begin{tikzpicture}[scale=1.05, every node/.style={font=\normalsize}]
\draw (0,0) rectangle (10,8);
\filldraw[fill=blue!30, fill opacity=0.35, draw=blue!60!black, thick]
(4,4.8) circle (2.4);
\filldraw[fill=red!30, fill opacity=0.35, draw=red!60!black, thick]
(6,4.8) circle (2.4);
\filldraw[fill=green!30, fill opacity=0.35, draw=green!50!black, thick]
(5,3) circle (2.4);

\node[anchor=north west] at (0.2,7.8) {$\powerset{B}$};
\draw[blue!60!black, very thick] (2.3,7.6) -- (2.8,7.6);
\node[anchor=west] at (2.9,7.6) {$b \in C$};
\draw[red!60!black, very thick] (4.8,7.6) -- (5.3,7.6);
\node[anchor=west] at (5.4,7.6) {$c \in C$};
\draw[green!50!black, very thick] (7.3,7.6) -- (7.8,7.6);
\node[anchor=west] at (7.9,7.6) {$d \in C$};

\node at (1,1) {$\emptyset$};
\node at (2.6,4.75) {$\set{b}$};
\node at (7.4,4.75) {$\set{c}$};
\node at (5,1.45) {$\set{d}$};
\node at (5,6.15) {$\set{b,c}$};
\node at (3.55,3) {$\set{b,d}$};
\node at (6.45,3) {$\set{c,d}$};
\node at (5,4.25) {$\set{b,c,d}$};
\end{tikzpicture}
\captionof{figure}{Potenssijoukon \(\powerset{B}\) Vennin diagrammi.
Ympyrät kuvaavat
ehtoja \(b \in C\), \(c \in C\) ja \(d \in C\), missä \(C \subseteq B\).}
\end{minipage}
\end{center}
\end{solution}

\begin{exercise}
Olkoon \(A\) ja \(B\) joukkoja. Todista \emph{absorbiolait}
\begin{enumerate}[label=\alph*), nosep]
\item \(A \cup (A \cap B) = A\),
\item \(A \cap (A \cup B) = A\).
\end{enumerate}
Havainnollista tuloksia Vennin diagrameilla.
\end{exercise}

\begin{solution}
\leavevmode\par
\begin{enumerate}[label=\alph*), nosep]
\item Olkoon \(C = A \cup (A \cap B)\). Näytetään ensin \(C \subseteq A\).
Olkoon \(x \in C\). Joukon yhdisteen määritelmän mukaan joko \(x \in
A\) \emph{tai} \(x \in A \cap B\).
Ensimäisessä tapauksessa \(x \in A\), ja toisessa tapauksessa
\(x \in A \cap B \subseteq A\), joten myös \(x \in A\). Siis \(C \subseteq A\).

Näytetään seuraavaksi \(A \subseteq C\). Olkoon \(x \in A\).
Koska \(x \in A\), niin \(x \in A \cup (A \cap B) = C\).
Siis \(A \subseteq C\).

Olemme osoittaneet \(C \subseteq A\) ja \(A \subseteq C\), joten
\(A = C = A \cup (A \cap B)\).

\item Olkoon \(C = A \cap (A \cup B)\). Näytetään ensin \(C \subseteq A\).
Olkoon \(x \in C\). Joukon leikkauksen määritelmän mukaan \(x \in A\) \emph{ja}
\(x \in A \cup B\). Siis jokaisesta \(x \in C\) seuraa \(x \in A\), joten
\(C \subseteq A\).

Näytetään seuraavaksi \(A \subseteq C\). Olkoon \(x \in A\). Lisäksi
\(x \in A\) merkitsee yhdisteen määritelmän mukaan, että \(x \in A \cup B\).
On siis voimassa \(x \in A\) ja \(x \in A \cup B\), joten
\(x \in A \cap (A \cup B) = C\). Voimme siis todeta \(A \subseteq C\).

Olemme osoittaneet \(C \subseteq A\) ja \(A \subseteq C\), joten
\(A = C = A \cap (A \cup B)\).
\end{enumerate}

\begin{center}
\begin{minipage}{\textwidth}
\centering
\begin{tikzpicture}[scale=0.9, every node/.style={font=\small}]
% First absorption law
\begin{scope}
\fill[blue!25] (1.4,1.2) circle (0.85);
\draw (0,0) rectangle (3.6,2.7);
\draw (1.4,1.2) circle (0.85);
\draw (2.2,1.2) circle (0.85);
\node at (0.2,2.48) {$X$};
\node at (1.4,2.4) {$A$};
\node at (2.2,2.4) {$B$};
\node[below] at (1.8,0) {$A \cup (A \cap B) = A$};
\end{scope}

% Second absorption law
\begin{scope}[xshift=5cm]
\fill[blue!25] (1.4,1.2) circle (0.85);
\draw (0,0) rectangle (3.6,2.7);
\draw (1.4,1.2) circle (0.85);
\draw (2.2,1.2) circle (0.85);
\node at (0.2,2.48) {$X$};
\node at (1.4,2.4) {$A$};
\node at (2.2,2.4) {$B$};
\node[below] at (1.8,0) {$A \cap (A \cup B) = A$};
\end{scope}
\end{tikzpicture}
\captionof{figure}{Absorbiolakien visualisoinnit Vennin diagrammeina. Molemmissa
tapauksissa tulokseksi jää joukko \(A\).}
\end{minipage}
\end{center}
\end{solution}

\begin{exercise}
Olkoon \(A_k = \left]0, \frac{1}{k}\right[, k \in \NN \setminus \set{0}\).
Määritä joukot
\begin{tasks}(4)
\task \(\displaystyle \bigcup_{k=1}^n A_k\),
\task \(\displaystyle \bigcap_{k=1}^n A_k\),
\task \(\displaystyle \bigcup_{k=1}^{\infty} A_k\),
\task \(\displaystyle \bigcap_{k=1}^{\infty} A_k\).
\end{tasks}
\end{exercise}

\begin{solution}
\leavevmode\par
\begin{enumerate}[label=\alph*), nosep]
\item Olkoon \(n \in \NN\). Kun \(1 \leq a < b \leq n\),
niin \(\frac{1}{b} < \frac{1}{a}\), joten \(A_b \subset A_a\). Siis
\[
A_n \subset A_{n - 1} \subset \cdots \subset A_2 \subset A_1.
\]
Koska \(A_1 = ]0, 1[\), saadaan
\[
\bigcup_{k=1}^n A_k = \left]0, 1\right[.
\]

\item Olkoon \(n \in \NN\). Jälleen, kun \(1 \leq a < b \leq n\),
niin \(\frac{1}{b} < \frac{1}{a}\), joten \(A_b \subset A_a\). Siis
\[
A_n \subset A_{n - 1} \subset \cdots \subset A_2 \subset A_1.
\]
Koska \(A_n = ]0, \frac{1}{n}[\), saadaan
\[
\bigcap_{k=1}^n A_k = \left]0, \frac{1}{n}\right[.
\]

\item Koska \(A_k \subseteq A_1\) kaikilla \(k \geq 1\), niin
\[
\bigcup_{k=1}^{\infty} A_k \subseteq A_1.
\]
Toisaalta, \(A_1\) on yksi yhdisteen joukoista, niin
\[
A_1 \subseteq \bigcup_{k=1}^{\infty} A_k.
\]
Siis
\[
\bigcup_{k=1}^{\infty} A_k = \left]0, 1\right[.
\]

\item Olkoon \(x > 0\). Valitaan \(n_0 \in \NN\) niin suuri, että
\(n_0 > \frac{1}{x}\). Tällöin \(\frac{1}{k} < n_0\). Joten
\(x \notin A_{n_0}\). Siis mikään \(x \in \RR\) ei kuulu jokaiseen
joukkoon \(A_k\), joten
\[
\bigcap_{k = 1}^{\infty} A_k = \emptyset.
\]
\end{enumerate}
\end{solution}

\begin{exercise}
Kirjoita määritelmän avulla inkluusio
\[
A \cup (B \cap C) \subseteq (A \cup B) \cap (A \cup C)
\]
implikaationa ja todista se vastatodistuksen avulla.
\end{exercise}

\begin{solution}
Inkluusion määritelmän mukaan väite on kaikilla \(x \in X\) pätevä implikaatio
\[
\begin{aligned}
x \in A \cup (B \cap C)
&\implies x \in (A \cup B) \cap (A \cup C).
\end{aligned}
\]
Joukko-operaatiot avaamalla sama implikaatio voidaan kirjoittaa
predikaattilogiikan kielellä muodossa
\[
\begin{aligned}
&x \in A \lor (x \in B \land x \in C) \\
&\qquad\implies
\bigl((x \in A \lor x \in B) \land (x \in A \lor x \in C)\bigr).
\end{aligned}
\]
Todistetaan nyt väite vastatodistuksella eli osoittamalla implikaation
kontrapositiivinen muoto. Olkoon \(x \in X\) mielivaltainen ja oletetaan, että
\[
x \notin (A \cup B) \cap (A \cup C).
\]
Tällöin \(x \notin A \cup B\) tai \(x \notin A \cup C\). Ensimmäisessä
tapauksessa \(x \notin A\) ja \(x \notin B\), joten
\(x \notin B \cap C\). Toisessa tapauksessa \(x \notin A\) ja
\(x \notin C\), joten myös tällöin \(x \notin B \cap C\). Kummassakin
tapauksessa siis
\[
x \notin A \cup (B \cap C).
\]
Näin on osoitettu kontrapositiivinen implikaatio
\[
\begin{aligned}
x \notin (A \cup B) \cap (A \cup C)
&\implies x \notin A \cup (B \cap C),
\end{aligned}
\]
joten alkuperäinen implikaatio pätee.
Koska \(x \in X\) oli mielivaltainen,
\[A \cup (B \cap C) \subseteq (A \cup B) \cap (A \cup C).\]
\end{solution}

\begin{exercise}
Olkoon \(A\), \(B\) ja \(C\) joukkoja. Todista karteesisen tulon osittelulaki
leikkauksen suhteen:
\[
(A \cap B) \times C = (A \times C) \cap (B \times C).
\]
Havainnollista tulosta joukoilla \(A = \set{1, 2, 3}\), \(B = \set{2, 3, 4}\)
ja \(C = \set{1, 2}\) määrittämällä ensin joukko \((A \cap B) \times C\)
ja sitten joukko \((A \times C) \cap (B \times C)\).
\end{exercise}

\begin{solution}
Asetetaan
\[
\begin{aligned}
D &= (A \cap B) \times C \\
&= \set{(x, y) \colon x \in A \cap B, y \in C}, \\
E &= (A \times C) \cap (B \times C) \\
&= \set{(x, y) \colon x \in A, y \in C} \cap \set{(x, y) \colon x \in
B, y \in C}.
\end{aligned}
\]
Ensiksi, olkoon \((x, y) \in D\) mielivaltainen. Sitten
\[
\begin{aligned}
(x, y) \in D &\implies x \in (A \cap B) \land y \in C \\
&\implies (x \in A \land x \in B) \land y \in C \\
&\implies (x, y) \in (A \times C) \land (x, y) \in (B \times C) \\
&\implies (x, y) \in (A \times C) \cap (B \times C) \\
&= E,
\end{aligned}
\]
joten \(D \subseteq E\).

Toiseksi, olkoon \((x, y) \in E\) mielivaltainen. Sitten
\[
\begin{aligned}
(x, y) \in E &\implies (x, y) \in (A \times C) \land (x, y) \in (B \times C) \\
&\implies (x \in A, y \in C) \land (x \in B, y \in C) \\
&\implies (x \in A \land x \in B) \land y \in C \\
&\implies x \in (A \cap B) \land y \in C \\
&\implies (x, y) \in (A \cap B) \times C \\
&= D,
\end{aligned}
\]
joten \(E \subseteq D\).

Olemme osoittaneet \(E \subseteq D\) sekä \(D \subseteq E\), joten
\[
(A \times C) \cap (B \times C) = E = D = (A \cap B) \times C.
\]
Havainnollistetaan vielä tulosta tehtävänannon joukoilla \(A\), \(B\) ja \(C\).
Saadaan
\[
\begin{aligned}
(A \cap B) \times C &= \set{(x, y) \colon x \in (\set{1,2,3} \cap
\set{2,3,4}), y \in \set{1, 2}} \\
&= \set{(x, y) \colon x \in \set{2,3}, y \in \set{1, 2}} \\
&= \set{(2, 1), (2, 2), (3, 1), (3, 2)},
\end{aligned}
\]
sekä
\[
\begin{aligned}
(A \times C) \cap (B \times C) &= \set{(x, y) \colon x \in
\set{1,2,3}, y \in \set{1,2}} \\
&\cap \set{(x, y) \colon x \in \set{2,3,4}, y \in \set{1,2}} \\
&= \set{(2,1),(2,2),(3,1),(3,2),(4,1),(4,2)} \cap \set{(2, 1), (2,
2), (3, 1), (3, 2)} \\
&= \set{(2,1),(2,2),(3,1),(3,2)}.
\end{aligned}
\]
Joukkojen alkiot ovat siis identtiset molemmissa tapauksissa, mikä
havainnollistaa todistettua yhtälöä.
\end{solution}

\end{document}
